\documentclass[12pt,a3paper,landscape]{exam}
\usepackage{mystyles}
\usepackage{longtable}
\usepackage{mhchem}

\geometry{margin=0.7cm}
\pagenumbering{arabic}

\author{Juliette Fenogli}
\date{}

\title{grille_JF1}
\date{}

\pagestyle{empty}
\begin{document}
\tito{\textsc{JF 1}}


\thispagestyle{empty}


\begin{longtable}{|p{6cm}|l|m{10.5cm}|l|m{10.5cm}|l|m{10.5cm}|}
\hline
Nom &&&&&&\\[30 pt]
\hline
\multicolumn{7}{|l|}{\textbf{I) SYNTHÈSE D’UN COMPOSÉ ODORANT NATUREL : L’ACÉTATE DE BENZYLE}}\\
\hline
Question 1 (estérification) &&&&&&\\[30 pt]
\hline
Pesées de liquide ou pipette, solvant éprouvette &&&&&&\\[30 pt]
\hline
Montage (tube du Dean-Stark rempli à affleurement avec du cyclohexane), agitation &&&&&&\\[30 pt]
\hline
Question 2 : APTS &&&&&&\\[30 pt]
\hline
Question 3 : Quantité d'eau attendue Dean-Stark &&&&&&\\[30 pt]
\hline
CCM, prép échantillons piluliers, principe, qualité du dépôt, introduction dans la cuve, révélation (principe), rapport frontal, pince &&&&&&\\[70 pt]
\hline
Gestes extractions l-l, lavages: agitation puis dégazage vers le fond de la hotte, ampoule non bouchée au repos  &&&&&&\\[30 pt]
\hline
Question 4 : extraction produit phase aqueuse, rendement d'extraction  &&&&&&\\[30 pt]
\hline
Question 5 : lavage basique pour neutraliser, dégagement \ce{CO2} &&&&&&\\[30 pt]
\hline
Dernier lavage à \ce{NaCl}(sat.) : relarguage et/ou préséchage &&&&&&\\[30 pt]
\hline
Séchage phase orga : agitation de l'erlen, quantité de \ce{MgSO4} à ajouter, pulvérulence &&&&&&\\[30 pt]
\hline
Montage de filtration,
%(coton ou papier, sous vide ou pas), rupture du vide si il y a, 
ballon fixé et \textbf{taré} &&&&&&\\[30 pt]
\hline
Question 6 : indice de réfraction, principe de la technique, dépôt du liquide entre les deux traits, utilisation de l'oculaire et des molettes, nettoyage &&&&&&\\[30 pt]
\hline
Question 7 : purification par distillation ou chromatographie sur colonne : principe &&&&&&\\[30 pt]
\hline
 Question 8 : rendement &&&&&&\\[30 pt]
\hline


\multicolumn{7}{|l|}{\textbf{II) DOSAGE DU CURCUMIN CONTENU DANS LA POUDRE DE CURCUMA}}\\
\hline
Dilution : 10 mL introduits précisément à la pipette jaugée, 1 mL non précis à l'éprouvette, remplissage aux deux tiers puis agitation, puis trait de jauge, boucher, homogénéiser &&&&&&\\[30 pt]
\hline
Gamme étalon réalisée à partir de B avec fioles de 10 mL et pipette graduée ou burette.  &&&&&&\\[30 pt]
\hline
Q1 : Choix $\lambda$ travail pour minimiser $\Delta A$  et/ou augmenter la sensibilité &&&&&&\\[30 pt]
\hline
Mesures d'absorbance, blanc avec le solvant seul &&&&&&\\[30 pt]
\hline
Droite d'étalonnage, régression, incertitude sur les paramètres &&&&&&\\[30 pt]
\hline
Question 2 : détermination $\epsilon$ \textbf{molaire} (unité), incertitude &&&&&&\\[30 pt]
\hline
Pesée \textit{environ précisément} de 0,5 g de curcuma &&&&&&\\[30 pt]
\hline
Dissolution du curcuma, erlenmeyer fixé &&&&&&\\[30 pt]
\hline
Filtration dans la fiole, fiole fixée &&&&&&\\[30 pt]
\hline
Rinçage du solide et finalisation de la fiole &&&&&&\\[30 pt]
\hline
Dilution : 1 mL de la solution A' pipette, 1 mL d'acide à l'éprouvette, remplissage aux deux tiers puis agitation, puis trait de jauge, boucher, homogénéiser &&&&&&\\[30 pt]
\hline
Question 3 : calcul pourcentage massique curcumin, incertitude sur pente et l &&&&&&\\[30 pt]
\hline

\multicolumn{7}{|l|}{\textbf{III) ANALYSE D’UNE BOISSON ANTI-FATIGUE DÉSALTÉRANTE}}\\
\hline
Question 1 : proton le plus acide, base conjuguée stabilisée par mésomérie &&&&&&\\[30pt]
\hline
Notion de dosage simultané des deux acides, dont les 3 acidités de l'acide citrique, 2 équations à considérer &&&&&&\\[40pt]
\hline
Détermination de la concentration d'acide ascorbique par réaction avec le diiode, équation équilibrée &&&&&&\\[40 pt]
\hline
Détermination de la concentration d'acide citrique par soustraction &&&&&&\\[30pt]
\hline
Dosage acido-basique par la solution de soude. 
Si suivi pH-métrique : choix des électrodes, étalonnage ou non du pH-mètre (pas nécéssaire pour repérer l'équivalence), electrodes fixées, rincées, épongées (avant et après mesure).
Si suivi colorimétrique : choix de l'indicateur coloré : le pH à l'équivalence doit être dans la zone de virage. Utilisation de dozzzaqueux possible. &&&&&&\\[30pt]
\hline
Détermination de l'équivalence.
Si suivi pH-métrique : pas de méthode des tangentes, méthode de la dérivée.
Si suivi colorimétrique : titrage à la goutte près, dès que la couleur est persistante, usage papier blanc dessous pour mieux repérer le changement de couleur. Premier titrage rapide, puis 2ème précis.  &&&&&&\\[30pt]
\hline
Dosage de l'acide ascorbique par le diiode. Suivi colorimétrique. Travail en milieu acide (pas dismutation) &&&&&&\\[30pt]
\hline
Détermination de la concentration de chacun des deux acides. &&&&&&\\[30pt]
\hline
Discussion sur les sources d'erreur : \newline Concentration des solutions titrantes (solutions fraîches car dismutation I2 \& soude carbonatée) \newline Volumes prise d'essai : verrerie \newline Volume équivalent (lecture burette, verrerie burette, méthode, volume d'une goutte). &&&&&&\\[30pt]
\hline

\multicolumn{7}{|l|}{\textbf{Evaluation transverse}}\\
\hline
Sécurité &&&&&&\\[30pt]
\hline
Aisance &&&&&&\\[30pt]
\hline
Dynamisme &&&&&&\\[30pt]
\hline
Traitement des déchets&&&&&&\\[30pt]
\hline
Vaisselle &&&&&&\\[30pt]
\hline

\end{longtable}


\end{document}     



